\documentclass[11pt,a4paper]{article}
\usepackage[T1]{fontenc}
\usepackage[utf8]{inputenc}
\usepackage{lmodern}
\usepackage[margin=25mm]{geometry}
\usepackage{amsmath,amssymb,booktabs,tabularx,xurl,xcolor,hyperref}
\hypersetup{colorlinks=true,linkcolor=blue!50!black,urlcolor=blue!50!black,
  pdftitle={Inferring Aggregation Algebra and Minimal Mode Summaries},pdfauthor={Mettapedia Working Notes}}
\setlength{\emergencystretch}{3em}
\setlength{\parindent}{0pt}
\setlength{\parskip}{5pt}
\newcommand{\code}[1]{\nolinkurl{#1}}
\title{Inferring Aggregation Algebra\\and Minimal Mode Summaries}
\author{Mettapedia Working Notes}
\date{31 October 2026}
\begin{document}
\maketitle

\begin{abstract}
Five Lean modules connect the sorted run-count fold used in label propagation
to histogram aggregation, implement a proof-producing classifier for additive
counter updates, and characterize the smallest summary for two-label mode under
future additions. Within a declared affine update language, admission is both
sound and complete for translation updates and invariant under semantically
equivalent spellings. The optimized runner preserves ordinary execution even
on refusal. A separate contextual theorem shows that the two-label optional
mode needs exactly a vote margin and an occupancy flag; fixing the empty
default to the smaller label removes even the flag. These are source-model and
information laws, not an arbitrary MeTTa algebra-discovery frontend or a native
performance claim.
\end{abstract}

\section{A bag function is not necessarily a homomorphism}

A pure fold after sorting natural labels is insensitive to their original
permutation. That makes its result a function of a bag. It does not make that
result mergeable. The mode counterexample is small:
\[
 \operatorname{mode}[1,1,2]=\operatorname{mode}[1,1,1]=1,
\]
but appending $[2,2]$ makes the first winner $2$ and leaves the second winner
$1$. Thus no binary operation on just the two partial winners computes every
combined winner. Lean proves this impossibility universally over the proposed
merge operation.

The correct decomposition is
\[
 \text{input collection}\ \xrightarrow{\ h\ }\ \text{mergeable summary}
 \ \xrightarrow{\ r\ }\ \text{answer},
 \qquad h(x\mathbin{+\!+}y)=h(x)\oplus h(y).
\]
For mode, $h$ can be a histogram and $r$ smallest-label argmax. The readout
need not itself preserve $\oplus$. A useful region certificate therefore
identifies the summary algebra and readout, as well as the allowed observations.
The free commutative monoid property determines a \emph{homomorphism} from its
singleton action; it does not synthesize a suitable lifting for every bag
function by itself.

This sits in the classical constructive-algorithmics tradition of list
homomorphisms. Gibbons' third homomorphism theorem derives a homomorphism when
the same list function has both leftwards and rightwards formulations
\cite{gibbons}. The classifier below instead uses an independently proved
matrix criterion for a specific update language. It is not a formalization
of that general theorem.

\section{Deriving the actual sorted run fold}

The inspected CDLP source sorts the collected labels, then folds a tuple
containing current label/count and best completed label/count. Equal labels
increase the current count. A new label finalizes the previous run, replacing
the best only when its count is \emph{strictly} larger. A final comparison
closes the last run. Empty vote lists are handled by the enclosing source,
which retains the vertex's old label.

\code{SortedRunMode} models those updates directly. Its invariant decomposes
the processed list into
\[
 \text{completed runs}\mathbin{+\!+}[c,c,\ldots,c],
\]
where the suffix length is the current count, every completed label is smaller
than $c$, and the saved best is already the smallest modal label of the
completed part. The proof of \code{append_run_best} establishes why a strict
count increase is the right update: a tie must keep the earlier, smaller label.
Induction through the fold proves the source result equals the independent
histogram mode specification on every sorted input.

The mathematical optional label encodes a sentinel versus a natural label;
its \code{none} does not mean an evaluator empty-answer event. The source's
empty-input bypass and old-label fallback are proved separately.

\code{ModeAggregation} independently implements a counter-update loop and
proves it computes list multiplicities. Histograms add pointwise across
concatenation. A lexicographic score $(\text{count},\text{reverse label})$
gives the greatest count and smallest label on ties, regardless of candidate
enumeration order. The sparse finite-support interpretation has the same
readout. Consequently,
\[
 \operatorname{sourceMode}(\operatorname{sort}(x))
   =\operatorname{readout}(\operatorname{histogram}(x)).
\]
Controls show that skipping sorting breaks the source run counter, changing
strict comparison to non-strict comparison breaks ties, and forgetting
multiplicity changes the winner.

\section{A classifier that constructs its proof}

\code{AdditiveFoldInference} provides a finite, pure update language for $n$
natural-valued registers. Expressions contain register references, input-only
weights, addition and multiplication by natural constants. Weights contain
constants, the input value, equality indicators, addition and constant scaling.
All expressions are total; effectful calls are not hidden inside the weight
syntax.

Normalization constructs a coefficient vector and a residual input expression:
\[
 \operatorname{eval}(e_i;s,x)=\sum_j A_{ij}s_j+w_i(x).
\]
The algorithm checks whether $A=I$. On acceptance it returns the extracted
weights and constructs a Lean certificate of
\[
 \operatorname{step}(s,x)=s+w(x).
\]
The theorem is not an input supplied by the caller. It follows from recursive
normalization and the computed finite matrix check.

The criterion is exact within this regime:
\[
 A=I\quad\Longleftrightarrow\quad
 \forall s,x:\ \operatorname{step}(s,x)
      =s+\operatorname{step}(0,x).
\]
Necessity follows by probing zero and each basis state. Therefore programs
with identical transition semantics inside this IR have identical admission
status, even when reassociation, zero-scaled registers or harmless extra terms
disguise their source form. A control with zero-scaled state and reordered
indicator terms is accepted. A decimal accumulator $s\gets10s+x$ is refused
and demonstrably order-sensitive. Refusing a program proves only that it is
outside this translation regime, not that no other algebra can optimize it.

For accepted programs, the derived summary is
\[
 H(x_1\cdots x_m)=\sum_{k=1}^m w(x_k),\qquad
 \operatorname{run}(s_0,x)=s_0+H(x).
\]
It is a bag homomorphism; any pure readout inherits permutation invariance.
The executable inferred runner chooses this summary when admitted and the
ordinary fold otherwise. \code{runInferred_exact} proves unconditional equality
with ordinary execution for every program of the IR, and the runner composes
over sequential input segments. Refusal is not an error or an empty result.

This is proof-producing inference for a declared family, not automatic
extraction from arbitrary MeTTa functions. The CDLP sorted-fold-to-histogram
translation is a separate proved transformation. Its discovery has not been
delegated to an unverified finite testing oracle.

\section{Grouped execution and label-domain admission}

The histogram update $s_i\gets s_i+[x=i]$ is accepted by the classifier for
every finite register domain. \code{GroupedModeRegion} extends the resulting
counter vector by zero outside that domain and proves exact readout when
every actual label lies inside it. An explicit negative control shows that
ignoring an out-of-domain label loses the answer.

The existing \code{JoinAggregation.groupFold} supplies one product-monoid fold
over group keys. Projection at each key equals folding its filtered vote list,
and partial group summaries merge pointwise. The joined theorem connects this
carrier to the sorted source mode, including ties and duplicate solutions.
It preserves the aggregate of the actual collected occurrences, without
assuming graph-incidence multiplicities from a source comment.

This licenses a mathematical count-plus-argmax region. Native tuple layout,
counter overflow, state revision, publication of the next label generation,
source return counts and concurrent readers are further refinement obligations.
The proof does not replace CDLP's outer iteration schedule.

\section{A smaller summary derived from future observations}

For two labels $\mathsf{false}<\mathsf{true}$, let
\[
 D(x)=\#\mathsf{true}(x)-\#\mathsf{false}(x),\qquad
 E(x)=[x\ne[]].
\]
The summary $(D,E)$ combines as
\[
 (d,e)\oplus(d',e')=(d+d',\ e\lor e').
\]
An occupied summary returns true exactly when $d>0$; an unoccupied summary
returns no modal label. The readout is proved equal to independent counting
and then connected to the original sorted mode via labels $0$ and $1$.
The reachable summaries are exactly those where unoccupied implies zero
margin; combining reachable summaries preserves this invariant.

The stronger result is continuation minimality. Define two inputs to be
equivalent when every possible future added suffix produces the same optional
modal answer. Then Lean proves
\[
 \bigl(\forall z,\ \operatorname{mode}(xz)=\operatorname{mode}(yz)\bigr)
 \quad\Longleftrightarrow\quad (D(x),E(x))=(D(y),E(y)).
\]
To distinguish unequal margins, add votes that bring one margin exactly to
zero and the other across the tie boundary. Every integer difference has an
explicit nonempty vote-list realization. Empty and balanced nonempty inputs
are separated by the empty future suffix. Every summary able to answer all
future additions must distinguish these classes.

Thus a list with one vote of each kind and a list with two votes of each kind
have the same future-addition behavior. This is not permission to forget counts
universally. Deleting one vote of each kind empties the first but leaves the second
nonempty. Full histograms have a separately proved one-occurrence erase law;
the additive merge theorem alone does not establish deletion maintenance.
The same summary cannot recover total vote count: the two example lists have
lengths two and four. A source count, evidence multiplicity or work-accounting
observer needs additional retained information.

Even occupancy is observer-relative. If the caller fixes the empty default
to false, the least label, both an empty and a balanced nonempty input return
false. In that profile,
\[
 \bigl(\forall z,\ \operatorname{defaultFalse}(xz)
                  =\operatorname{defaultFalse}(yz)\bigr)
 \quad\Longleftrightarrow\quad D(x)=D(y).
\]
A true empty default separates the previously merged empty/balanced cases.
The known caller context therefore changes which information can be erased.

Finally, no finite-state summary can answer all future additions on unbounded
inputs: arbitrarily many integer margins have distinguishing contexts. A fixed
number of integer fields is not a fixed number of bits. The result is an exact
semantic information boundary, not a universal constant-space implementation.

\section{Scope and next refinements}

The implemented classification is complete for additive translation semantics
inside the affine IR and incomplete as a classifier of arbitrary programs.
The minimality theorem concerns two-label mode, not arbitrary many-label modes.
No theorem states that the fastest implementation follows automatically from
a semantic description. Admission supplies permission; physical cost models
and measurements must choose an implementation.

The next refinements are a source/profile extractor with revision and effect
certificates; broader summary inference for run folds and other algebraic
families; larger-domain continuation quotients; and native grouped execution
with accounting, lifetime and cost correspondence. Metamorphic tests should
vary equivalent source spellings and deliberately change effects, multiplicity,
tie policy, label domains and observable boundaries.

The modules under \code{Mettapedia/Machines/} are \code{ModeAggregation},
\code{SortedRunMode}, \code{AdditiveFoldInference}, \code{GroupedModeRegion},
and \code{BinaryModeInformation}. \code{AggregationInferenceAudit} inventories
their checked declaration bodies and transitive primitive axioms.

\begin{thebibliography}{9}
\bibitem{gibbons} Jeremy Gibbons, ``The Third Homomorphism Theorem,''
\emph{Journal of Functional Programming} 6(4), 657--665, 1996.
\href{https://doi.org/10.1017/S0956796800001908}{doi:10.1017/S0956796800001908}.
\end{thebibliography}
\end{document}
